% !TeX root = ../../Book2.tex
% 纲要标题：## 附录 B　分层空间、局部拓扑与交叉同调
% 纲要位置：几何学计划纲要.md，第 1941 行。
% 正文状态：未撰写；本文件目前是框架。

% 纲要内容开始
% 
% > 基础部分按“分层—控制—局部拓扑—允许链”深入。Whitney 和 Thom 型条件各有定义；一般层论对偶和中间延拓不冒充本附录已完成的全套理论。
% 
% * **分层的几何条件**
%   * 局部有限分层与边界条件
%   * Whitney (a)、(b) 的切空间极限
%   * 横截切片与链接
%   * 代数、解析和光滑分层的区别
% 
% * **控制数据与局部平凡性**
%   * 管状控制数据及相容性
%   * 分层向量场和积分
%   * 第一同位引理的选定固有版本
%   * 缺少控制或固有性时的失败
% 
% * **锥、链接与允许链**
%   * 局部锥模型与普通对偶的失败
%   * 允许函数（perversity）的维数限制与允许链
%   * 锥公式的代表性计算
%   * 交叉同调的对偶定理：精确陈述及证明输入
% 
% * **层方法与消失圈**
%   * 可构造层和局部系统
%   * 交叉链层与中间延拓的动机
%   * 平滑族与退化族的附近圈
%   * 结点退化的消失圈计算
% 
% ---
% 
% 纲要内容结束

\chapter{分层空间、局部拓扑与交叉同调}
\label{g2:app:B}
\BookChapterNumber{g2:app:B}

待撰写

\section{分层的几何条件}
\label{g2:sec:B01}
待撰写

\subsection{局部有限分层与边界条件}
\label{g2:subsec:B01:01}
待撰写

\subsection{Whitney (a)、(b) 的切空间极限}
\label{g2:subsec:B01:02}
待撰写

\subsection{横截切片与链接}
\label{g2:subsec:B01:03}
待撰写

\subsection{代数、解析和光滑分层的区别}
\label{g2:subsec:B01:04}
待撰写

\section{控制数据与局部平凡性}
\label{g2:sec:B02}
待撰写

\subsection{管状控制数据及相容性}
\label{g2:subsec:B02:01}
待撰写

\subsection{分层向量场和积分}
\label{g2:subsec:B02:02}
待撰写

\subsection{第一同位引理的选定固有版本}
\label{g2:subsec:B02:03}
待撰写

\subsection{缺少控制或固有性时的失败}
\label{g2:subsec:B02:04}
待撰写

\section{锥、链接与允许链}
\label{g2:sec:B03}
待撰写

\subsection{局部锥模型与普通对偶的失败}
\label{g2:subsec:B03:01}
待撰写

\subsection{允许函数（perversity）的维数限制与允许链}
\label{g2:subsec:B03:02}
待撰写

\subsection{锥公式的代表性计算}
\label{g2:subsec:B03:03}
待撰写

\subsection{交叉同调的对偶定理：精确陈述及证明输入}
\label{g2:subsec:B03:04}
待撰写

\section{层方法与消失圈}
\label{g2:sec:B04}
待撰写

\subsection{可构造层和局部系统}
\label{g2:subsec:B04:01}
待撰写

\subsection{交叉链层与中间延拓的动机}
\label{g2:subsec:B04:02}
待撰写

\subsection{平滑族与退化族的附近圈}
\label{g2:subsec:B04:03}
待撰写

\subsection{结点退化的消失圈计算}
\label{g2:subsec:B04:04}
待撰写
